\cleardoublepage
\thispagestyle{plain}
\addcontentsline{toc}{chapter}{ÖZET}
\begingroup
\singlespacing

\vspace*{72pt}
\begin{center}
  \bfseries\MakeUppercase{\TezBasligi}\\[18pt]
  ÖZET
\end{center}

% Not: Nihai benzetim/deney sonuçları kilitlendiğinde sayısal değerler gözden geçirilecektir.
Enerji depolama sistemleri, yenilenebilir kaynakların şebekeye entegrasyonunda üretim ile tüketimin dengelenmesini sağlayan temel yapılardan biridir ve sistem başarısı yalnız batarya kapasitesiyle değil, güç dönüşüm arayüzünün verimi, güvenilirliği ve maliyetiyle de belirlenir. Şebekeye bağlı üç seviyeli T-tipi evirici, anma gücünde yüksek verim ve güç yoğunluğu sunmasına karşılık, çalışma süresinin önemli bir bölümünü oluşturan düşük güç bölgesinde farklı bir sorunla karşılaşır: sıfır gerilim durumunu üreten nötr anahtar kolu, endüktör akımı sıfıra ulaştıktan sonra iletimde tutulduğunda senkron alçaltıcı dönüştürücü benzeri bir ters akım dolaşımı oluşturur. Yararlı güç taşımayan bu dolaşım, RMS akımı ve iletim kayıplarını artırarak hafif yük verimini düşürür.

Bu tezde söz konusu mekanizma; anahtarlama durumları, akım yolları ve senkron-buck eşdeğeri üzerinden çözümlenmiş ve üç özgün katkı tek bir araştırma zinciri altında birleştirilmiştir. Birinci katkı, ters akımı sürdüren nötr anahtar komutunu düşük güç bölgesinde koşula bağlı olarak bastıran yeni bir anahtarlama yöntemidir. Yöntem benzetim ve prototip kayıtlarında verimi artırırken akımı süreksiz iletime zorlamakta ve üçüncü harmonik ağırlıklı belirgin bir bozulma oluşturmaktadır. İkinci katkı, bu harmonik cezayı azaltmak için şebeke faz bilgisinden üretilen Sin\textsuperscript{3} tabanlı ileri besleme kompanzasyonudur. Terim, hata sinyalinden türetilen bağımsız bir kontrolcü değil, mevcut akım döngüsünün modülasyon toplamına eklenen bir bileşendir; kapalı çevrim modeli, duyarlılık fonksiyonları ve kararlılık marjı incelemesiyle birlikte ele alınmıştır. $\sin^3\theta=[3\sin\theta-\sin(3\theta)]/4$ özdeşliğine dayanan parçalı DCM--CCM modeliyle yapılan katsayı taraması, kompanzasyon katsayılarının çalışma noktasına bağlı olduğunu göstermiş ve doğrulanmış parametre seti için üçüncü harmonik oranını analitik modelde yüzde 39,96'dan yüzde 0,15'e indiren düzeltilmiş katsayı adayları üretmiştir. Üçüncü katkı, aynı T-tipi yarı iletkenlerin evirici, Vienna PFC ve aktif batarya yönetimi işlevlerinde yeniden görevlendirilmesine dayanan GT-tipi yeniden yapılandırılabilir enerji depolama mimarisidir; çalışma kipleri, güvenli geçiş koşulları ve doğrulama matrisi tanımlanmıştır.

Benzetim kayıtlarında nötr anahtar susturma çıkış akımı THD değerini yüzde 2,14'ten yüzde 46,35'e yükseltmiş; Sin\textsuperscript{3} kompanzasyonu aynı değeri yüzde 7,79'a indirmiştir. Prototip deneyleri düşük güçte verim artışı eğilimini desteklemektedir. Benzetim ve deney sonuçları ayrı raporlanmış, her iddianın kanıt sınırı açıkça belirtilmiş ve tamamlanması gereken doğrulamalar metin içinde işaretlenmiştir.

\textbf{Anahtar Kelimeler:} T-tipi evirici, enerji depolama sistemi, düşük güç, nötr anahtar susturma, ileri besleme kompanzasyonu, üçüncü harmonik, yeniden yapılandırılabilir mimari.
\endgroup
